\section{The scalar in a rational graph term}

% BEGIN E76Q1
\paragraph{E76Q1. Transport of the two elliptic variables and the Gaussian.}
Take the graph matrix $g=i_\gamma$ of O74E1, so the frame pair
$(x,y)$ moves to $(x,\gamma y)$. Fix a CM place and write the
blocks of $g$ there as $A,B,C,D$. Put
\[
 Q=CZ+D,\qquad Q'=ZC^*+D^*,\qquad
 W=(AZ+B)Q^{-1},\quad Y=\operatorname{Im}_H Z .
\]
Here $Z$ is the doubled period. Write $L=\mathrm{aLin}$ and
$H=\mathrm{aAnti}$ for the two invertible frame matrices used to
define this period, with $G$ the embedded Hermitian Gram matrix.
Let $P$ select the first signature block and set
$p=L^{-1}PL$, $p'=H^{-1}PH$. In the following computation $u$
is the image of $\gamma^{-1}$ at this place. Projection onto the
two blocks in the definitions of $L,H$ gives the moved matrices
$LM,HM'$, where
\[
 M=pu+(1-p),\qquad M'=p'+(1-p')\bar u.
\]
Indeed the selected rows of $L$ transform on the right by $u$,
and the complementary rows are fixed; for $H$ the complementary
rows transform by $\bar u$. This uses $\gamma^tG\bar\gamma=G$.
The mixed matrix $HG^{-1}L^t$ commutes with $P$: its off-block
entries are those of the identity between a frame matrix and its
inverse, in opposite blocks. This gives $p^tG=Gp'$.

Substitute $Z=c((1-2p)^tG)$, $c=\tau(\kappa_0)/4$, into the blocks
of $i_\gamma=T^{-1}\operatorname{diag}(1,u^t)T$ from O74E1.
Multiplying out and using $p^2=p$ gives
\[
 GQ=M^tG,\qquad Q'G=GM',\qquad
 AZ+B=c\big((1-2p)M\big)^tG.
\]
These identities give the doubled period of the moved frames as
$W$ and carry an elliptic column $v$ to $(Q')^{-1}v$ and an
independent row $v'$ to $v'Q^{-1}$. All inverses exist on the
frame domain. In fact $Q^*Y_W Q=Y$ for $Y_W=\operatorname{Im}_H W$,
by the split unitary block identities, and $Y$ is positive
definite. Invertibility of $Q$ also follows directly from
$GQ=M^tG$ and invertibility of the moved frame matrix $LM$.

The Gaussian contribution can be verified without a choice of
square root. Because $CD^*=DC^*$, one has $QC^*=CQ'$ and
$Q^*=Q'-2iYC^*$. It follows that
\[
 Q^{-1}Y_W^{-1}(Q')^{-1}+2iQ^{-1}C=Y^{-1}.
\]
Pairing with $v',v$ proves that the two exponentials in the moved
Gaussian and the factor of automorphy give the original exponential.
The determinant factors satisfy
$\det M=\det Q$, $\det M'=\det Q'$ and
$\det Q=\det(g)\det Q'$. For the last equality,
$Q'(A-WC)=1$ by the unitary identities, and block elimination in
$g\left(\begin{smallmatrix}1&Z\\0&1\end{smallmatrix}\right)$
gives $\det g=\det(A-WC)\det Q$. Thus for $k=n/2$ the moved frame
factor times $\det(Q)^n$ leaves $\det(g)^k$ times the original
frame factor. Taking the product over the CM type gives exactly
\[
 k_0(x,\gamma y;z,w)\,\mathrm{hFac}(g;p_{x,y})=
 \left(\prod_{\tau\in\Psi}\tau(\det i_\gamma)^{n/2}\right)
                 k_0(x,y;z,w).
\]
Here $p_{x,y}$ denotes the doubled period and elliptic variables.
The displayed scalar is nonzero and independent of the frames and
elliptic arguments.
% END E76Q1

% BEGIN E76Q2
\paragraph{E76Q2. The unitary move on the original finite test.}
For an ordinary unitary move $\gamma$, identify its action on the
rational period lattice in the chosen basis. Preservation of the
trace pairing makes that rational matrix symplectic. Its finite
Weil operator, built from coordinate substitution, multiplication
by a unitary quadratic character and Fourier transformation, preserves
the finite Haar inner product. On the ordinary Siegel theta function
it gives the theta change of variables (Poisson summation for the
Fourier generator). In the framed reading, the inverse period matrix
changes by the inverse transpose of the Siegel denominator. Its
determinant cancels the weight $n/2$ factor; the quadratic change
of the exponent cancels the symplectic exponential, by the
imaginary-part identity for the Siegel period. Hence with the
native operator $w_G$,
$\Theta_A(w_G(\gamma)f;\gamma y)=\Theta_A(f;y)$ in all elliptic
variables. The operator so specified is unique on finite tests:
division by the nonzero frame prefactor and Fourier coefficient
extraction at a fixed period recover each value of $f$. The same
uniqueness gives its Heisenberg covariance by applying the two
orders of operations to the translated theta formula; the alternating
trace phase is preserved by $\gamma$.

% END E76Q2

% BEGIN E76Q3
\paragraph{E76Q3. Evaluation on the graph of the unitary move.}
Apply E76Q2 and the doubled product formula to the transformed
frames in E76Q1. The two doubled theta functions for
$\omega(i_\gamma)\mathcal D(f,h)$ and
$\mathcal D(f,w_G(\gamma)h)$ then differ by exactly the determinant
scalar of E76Q1. At the moved period every pair of elliptic
variables is attained since the two frame matrices are invertible.
Taking Fourier coefficients recovers the same equality of tests.
Evaluation at zero and division by the Eisenstein automorphy
factor cancel that scalar, leaving
$k_0(x,\gamma y)\mathcal D(f,w_G(\gamma)h)(0)$. By O74E2 this is
$k_0(x,\gamma y)a\langle f,w_G(\gamma)h\rangle$, with its single
nonzero constant $a$. For the system tests the support and
covariance of L76F1--L76F2 give exactly the graph coefficient of
O74E3, including the conjugate of the level multiplier. On a
nongraph isotropic space its first-block direction gives a pure
modulation after the symplectic change of coordinates. Apply the
annihilator of L76F3 and evaluate at zero as in O74E3. Thus the
entire graph and singular-term calculation uses the same finite
tests as the geometric readings.
% END E76Q3
